\documentclass{article}
\usepackage[utf8]{inputenc}
\usepackage[T1]{fontenc}
\usepackage{geometry}
\usepackage{graphicx}
\usepackage{float}
\usepackage{hyperref}
\usepackage[english,indonesian]{babel}

\title{Optimasi Pemrosesan Template LaTeX Menggunakan Pendekatan Berbasis AI}
\author{Academic Support Agent}
\date{\today}

\begin{document}
\maketitle

\begin{abstract}

\end{abstract}

\chapter{PENDAHULUAN}


\chapter{TINJAUAN PUSTAKA}


% Original content (commented out):
% \documentclass[a4paper]{article}

\usepackage[english]{babel}
\usepackage[utf8x]{inputenc}
\usepackage{amsmath}
\usepackage{graphicx}
\usepackage[colorinlistoftodos]{todonotes}

\title{2.2 Formal Write Up}
\author{Jack Bloomer}

\begin{document}
\maketitle
\section*{{Problem 2.2.10}}

\textit
{In this problem, you will create a model of a crane population that includes “hacking” chicks, which means to hatch chicks in captivity and release them into the wild. 
Assume that your crane population has a constant birth rate per year of 2\% of the 
population and a constant death rate of 2.5\%. Assume also that 5 chicks are “hacked” 
each year; i.e. assume 5 hatched chicks are released into the wild and absorbed into the 
population.}
\newline
\newline
\textbf{1. Draw a compartmental diagram for your model.}
\newline
\begin{figure}[h]
\includegraphics{cd.png} 
\end{figure}
\newline
\newline
\textbf{2. Use the compartmental diagram to create your model.}
\newline
\newline
Looking at this compartmental diagram, we can see that the model for this problem will 
determine each year’s population by having the birth rate added to the population, 5 hacked 
chicks added into the population, and the death rate taken out of the population. This model 
can be represented as such:
\[P(t+1)=0.995P(t)+5\]
with P representing the population and t representing time in years.
\newline
\newline
\newline
\newline
\newline
\textbf{3.What are the fixed points of your model? Are they stable or unstable? How can you
interpret this physically?}
\newline
\newline
To find the fixed point of this model, we can find when $f(\bar{a})=\bar{a}$.
$$P(t+1)=0.995P(t)+5$$
$$\bar{a}=0.995\bar{a}+5$$
$$0.005\bar{a}=5$$
$$\bar{a}=1000$$
\newline
To find out if this fixed point is stable or unstable, we need to find out if 
$|f'(\bar{a})|$ is less than or greater than 1.
$$P(t+1)=0.995P(t)+5$$
$$f(\bar{a})=0.995\bar{a}+5$$
$$f'(\bar{a})=0.995$$
\newline
Since $f'(\bar{a})$ is less than one, this tells us that the fixed point is stable, or in 
other words, an attracting fixed point. If we were to interpret our fixed point physically, 
we would say that once the population reaches 1,000, it will no longer increase nor will it 
decrease because the birth and death rate will result in a population of 995 plus the 5 
hacked chicks every year from there on out.

\end{document}

\end{document}